\documentclass[12pt]{article}
\usepackage{geometry}
\usepackage{setspace}
\usepackage{graphicx}
\usepackage{hyperref}
\usepackage{amsmath}
\usepackage{booktabs}
\usepackage{array}

\geometry{a4paper, margin=1in}
\onehalfspacing

\title{Baselines Phase Report \\ \large Reproducibility of the GALOIS System}
\author{
    Your Name 1 \\
    Your Name 2 \\
    Your Name 3
}
\date{19 November 2025}

\begin{document}

\maketitle

\begin{abstract}
This project is part of the reproduction of the GALOIS system, which aims to improve SQL query execution over Large Language Models (LLMs). Since LLMs are increasingly used for structured information retrieval, it is important to understand how well they behave when treated as a ``virtual database.''

In this phase, we implemented and evaluated several baseline methods: direct SQL prompting, natural language prompting, and few-shot prompting. These baselines were executed across three different LLM providers: Google Gemini, IBM WatsonX, and Azure OpenAI. The purpose of this report is to describe our methodology, the experimental setup, the results, and how this work prepares the ground for the next phase, where we will implement the optimized strategies introduced in the GALOIS framework.
\end{abstract}

\section{Introduction}

Large Language Models contain a significant amount of factual information encoded in their parameters. Several recent works, including the GALOIS paper, examine whether this information can be accessed through SQL queries instead of natural language. In practice, however, LLMs often output incomplete or partially incorrect tables.

Our long-term objective is to fully reproduce the GALOIS system, including its logical and physical query optimization strategies. The first step toward this objective is to implement the baseline systems used for comparison in the original paper. Furthermore, to observe the model-specific differences, we executed each baseline on three LLM providers:

\begin{itemize}
    \item Google Gemini
    \item IBM WatsonX
    \item Azure OpenAI (GPT models)
\end{itemize}

This gives us a broader understanding of the behavior of different LLM families under the same tasks.

\section{Experimental Setup}

To ensure reproducibility, we describe our setup in detail.

\subsection{Software}
\begin{itemize}
    \item Python 3.10
    \item APIs/SDKs for:
    \begin{itemize}
        \item Google Gemini
        \item IBM WatsonX
        \item Azure OpenAI
    \end{itemize}
    \item requests library for handling HTTP requests
    \item Pandas, NumPy, DuckDB
    \item JSON parsing utilities like orjson
    \end{itemize}

\subsection{Models Used}
\begin{itemize}
    \item Gemini 2.5 Flash/Pro
    \item WatsonX Granite model 3.3.18
    \item Azure OpenAI GPT-4.0
\end{itemize}

\subsection{Data and Queries}

We used:
\begin{itemize}
    \item The NL and SQL query files from the given datasets
    \item Ground truth that can be generated by \texttt{generate\_ground\_truth.py}
    \item Table schemas defined in the GALOIS paper for comparing to our results
\end{itemize}

\subsection{Evaluation Metrics}
We computed:
\begin{itemize}
    \item F1-Cell
    \item Cardinality
    \item Tuple Constraint
    \item Average Score (the Average value of three metrics)
    \item Token usage (when available by model)
    \item Average Time
\end{itemize}

Additionally, we measured absolute and relative differences compared to the published numbers in paper.

\section{Baseline Methods}

The same three baselines were applied to all three LLM providers.

\subsection{Direct SQL Prompting}

The SQL query is sent directly to the LLM.  
We instruct the model to return a JSON table that follows the expected schema.  
This represents the simplest approach and is one of the baselines used in the original GALOIS paper.

\subsection{NL Baseline (Baseline 2)}

In this part of the project we implemented the ``Natural Language Baseline'' which corresponds to Baseline~1. The goal of this baseline is to evaluate how well an LLM answers a query when it is a human-readable question. This gives us a reference point for understanding how much structure SQL adds compared to plain natural language.

\subsection{Code Structure}

The NL baseline is implemented in the file \texttt{nl\_baseline.py}~\cite{nl_baseline}.  
The script is responsible for:
\begin{itemize}
    \item discovering datasets,
    \item reading all NL queries for each dataset,
    \item Using each NL query with a prompt template,
    \item calling the selected LLM provider,
    \item extracting the JSON output in full fromat, 
    \item saving the result for evaluation.
\end{itemize}

We supported three LLM providers:
\begin{itemize}
    \item \textbf{Azure OpenAI (GPT models)},
    \item \textbf{Google Gemini},
    \item \textbf{IBM WatsonX}.
\end{itemize}

At runtime, the provider is selected through a command-line flag 
(\verb|--provider azure|, \verb|--provider gemini|, or \verb|--provider watsonx|),
and the script loads the corresponding client from 
\texttt{src/llm/azure\_client.py}, \texttt{gemini\_client.py}, or 
\texttt{watsonx\_client.py}.
If no provider is specified, the system defaults to \textbf{WatsonX}.
Likewise, if no dataset is specified, the default behavior is to use 
\textbf{all available datasets}.

\subsection{Natural Language Prompt Construction}

Each natural-language query (NL query) is converted into a prompt using the internal function \verb|_build_prompt(nl_query: str) -> str| .
The NL baseline does not attempt to rewrite or simplify the query.  
Instead, it takes the raw natural-language query and inserts it into a fixed instruction template.  
The prompt used is:

\begin{verbatim}
Answer the following query by giving a JSON array.
Make sure the output is valid JSON and uses consistent object keys.
No additional text is needed.

Query:
<the original NL query, unchanged>
\end{verbatim}

All LLM providers---Azure, Gemini, and WatsonX---therefore receive the same prompt format, ensuring that model differences arise from the LLMs themselves rather than prompt design.

\subsection{Dataset Loading and Query Enumeration}

Datasets are discovered using the helper function 
\texttt{discover\_datasets()} from \texttt{common.py}.  
For each dataset directory, the natural-language queries are loaded via
\texttt{read\_queries()}, which returns a list of NL query strings.  
The baseline code does not modify query structure; it simply loads all files named
\texttt{queries\_*.txt} in the dataset directory and extracts their contents.

\subsection{LLM Execution}

For each constructed prompt, the script calls the selected model backend through a unified interface.  
The backend is loaded dynamically using \verb|_load_provider(name) -> generate_text|


This function returns the appropriate \verb|generate_text(prompt)| method from: \texttt{azure\_client.py}, \texttt{watsonx\_client.py} and \texttt {gemini\_client.py}.

Each provider function returns either:
\begin{itemize}
    \item the generated text, or
    \item a tuple \verb|(text, tokens)| if token information is available.
\end{itemize}

Execution time is recorded for every query so that our evaluation can compare accuracy, response structure, and runtime.

\subsection{JSON Extraction and Error Handling}

Because LLM outputs may contain text surrounding the JSON, the baseline implements a lightweight JSON recovery mechanism:
\begin{itemize}
    \item try \verb|json.loads()| directly,
    \item if that fails, extract the text between the first ``['' and the last ``]'',
    \item if both fail, wrap the entire output into an error object.
\end{itemize}

All results are normalized into a standard structure (Full format):

\begin{verbatim}
{
  "result_set": [...],
  "tokens": <token_count>,
  "time": <execution_time>
}
\end{verbatim}

The output for each query is saved to \verb|nl_baseline_results/<DATASET>/queryX.json|.
Ensuring consistent formatting for future evaluation.

\subsection{Purpose of This Baseline}

The NL baseline establishes how LLMs behave when provided only natural-language instructions, without the structural information present in SQL or logical query plans.

This creates a ``worst-case'' reference point that motivates the design of the more structured GALOIS pipeline, where logical and physical plans reduce ambiguity and improve LLM determinism.


\subsection{Few-Shot Prompting}

We provide a few short examples of SQL queries paired with correct JSON tables.  
The idea is that the model will follow the pattern more consistently.

\section{Results}

\textit{(Your results tables go here. This section remains descriptive so that you can include your own measured values.)}

In general, the following trends were observed:

\subsection*{Natural Language Baseline}
\begin{itemize}
    \item Lowest recall across all three models
    \item WatsonX sometimes misinterpreted numeric filters
    \item Gemini often produced very short lists
    \item Azure generally produced the cleanest JSON
\end{itemize}

\subsection*{Direct SQL Baseline}
\begin{itemize}
    \item More structured than NL, but still incomplete
    \item Gemini sometimes hallucinated numbers
    \item WatsonX gave stable formatting but weak recall
    \item Azure was closest to ground truth but still missed data
\end{itemize}

\subsection*{Few-Shot Baseline}
\begin{itemize}
    \item Formatting improved significantly
    \item Azure benefited the most
    \item Gemini sometimes ignored examples
    \item WatsonX improved consistency but recall remained low
\end{itemize}

Overall:
\begin{itemize}
    \item Azure performed best in structure and precision
    \item Gemini had unstable recall depending on the query
    \item WatsonX was consistent but conservative (often missing tuples)
\end{itemize}

\section{Discussion}

Our experiments confirm that even strong LLMs struggle with structured SQL-like data extraction.  
Although the models differ, all of them had trouble achieving full recall, especially for queries requiring multiple conditions or less common entities.

Common issues included:
\begin{itemize}
    \item Invalid JSON (especially with Gemini and WatsonX)
    \item Hallucinated values (mostly Gemini)
    \item Missing entries (all three models)
    \item Strong sensitivity to prompt phrasing
\end{itemize}

These findings align with the motivation behind the GALOIS approach.

\section{Impact on Next Phase and Future Steps}

This baseline work directly influences the next part of the project.  
Key takeaways:

\begin{itemize}
    \item Simple prompting is unreliable, supporting the need for structured logical plans.
    \item The models frequently output only the ``easiest'' tuples, highlighting the importance of the Key-Scan operator.
    \item Behavior differs across providers, so our implementation should remain model-agnostic.
\end{itemize}

\subsection*{Next Phase Plan (Due 23 December)}
\begin{itemize}
    \item Implement logical plans (no pushdown, selective pushdown, full pushdown)
    \item Implement physical operators:
    \begin{itemize}
        \item Table-Scan
        \item Key-Scan
    \end{itemize}
    \item Add confidence-based plan selection
    \item Run all experiments again across the three model providers
    \item Compare against the baselines from this report
\end{itemize}

\section{Conclusion}

By testing Gemini, WatsonX, and Azure OpenAI with direct SQL, natural language, and few-shot prompting, we showed that LLMs still struggle with complete structured data extraction. Formatting quality varies between models, and all three miss relevant tuples in many queries.

This phase provides the necessary reference values for the next stage of the project, where we will reproduce the optimized strategies introduced by GALOIS and evaluate whether they significantly improve result quality.

\end{document}
